\chapter{连续模型与偏微分方程理论基础}
\label{ch:continuous-pde}

本章建立数值 PDE 的连续数学入口。计算科学中的离散方法并不是直接处理物理对象，而是首先把连续模型表示为微分算子、边界条件和初始条件组成的问题。

核心形式为
\begin{equation}
\mathcal{L}(u)=f.
\end{equation}

后续有限差分、有限体积和离散算子构造，都是建立离散算子
\begin{equation}
\mathcal{L}_h(u_h)=f_h
\end{equation}
的过程。

\section{连续模型与算子表示}

\section{偏微分方程的适定性}

\section{椭圆型方程}

\section{抛物型方程}

\section{双曲型方程}

\section{不同方程类型的数值特征}

\section{工程模拟中的耦合数学模型}
